\documentclass[12pt]{article}
\usepackage[utf8]{inputenc}
\usepackage[T1]{fontenc}
\usepackage{hyperref}
\title{Yote Language Resources for K-12 Curriculum Development}
\author{Compiled by Claude Sonnet 5 (claude-sonnet-5) for Andr\'as Kornai}
\date{}
\begin{document}
\maketitle

\section{Introduction}
Yote, also called Wab, is an Austronesian (Oceanic) language of Madang Province, Papua New Guinea (ISO 639-3 wab; Glottocode wabb1237). It has about 120 speakers in the villages of Wab and Saui. Little exists for it, and the search found only reference and archival listings. It should not be confused with Lote, an unrelated language of New Britain.

\section{Primary Educational Resources}
None found.

\section{Linguistic and Reference Materials}
\begin{itemize}
\item Yote language (Wikipedia): \url{https://en.wikipedia.org/wiki/Yote_language}
\item Wab language (Wikipedia): \url{https://en.wikipedia.org/wiki/Wab_language}
\item Numeral data for Yote (Max Planck database), which notes that the data is incomplete: \url{https://lingweb.eva.mpg.de/channumerals/Yote.htm}
\item Verbix endangered-languages page: \url{https://docs.verbix.com/EndangeredLanguages/Wab}
\end{itemize}

\section{Religious and Community Resources}
Joshua Project reports that Bible translation has started, with no audio Bible: \url{https://joshuaproject.net/languages/wab}

\section{Digital and Audio Resources}
Global Recordings Network has a listing for the language: \url{https://globalrecordings.net/en/language/wab}. I did not verify what it contains.

\section{Curriculum Development Context}
No school use is known. No dictionary, grammar or orthography was found. Any curriculum would have to start with community-led documentation.
\end{document}